% ============================ CONCLUSIONES ===========================
\section{Conclusiones}

\paragraph{1. La arquitectura desacoplada cumplió su propósito.} Separar el almacenamiento (Cloud Storage) del cómputo (clústeres efímeros de Dataproc) permitió destruir y reprovisionar dos topologías distintas sobre el mismo conjunto de datos sin ninguna pérdida ni recarga de información. Este desacoplamiento no es un detalle de implementación: es lo que hizo materialmente posible el experimento comparativo entre las Arquitecturas A y B que constituye el núcleo de este informe.

\paragraph{2. Para 500,000 registros, la herramienta más rápida fue la que no distribuye nada.} Polars completó las diez consultas en \textbf{0.567 segundos}, frente a 21.45\,s de Modin, 203.18\,s de Dask y 376.97\,s de Apache Spark. Spark resultó \textbf{665 veces más lento} que Polars ejecutando exactamente la misma lógica y produciendo resultados idénticos. La conclusión no es que Spark sea una mala herramienta, sino que 667.7 MiB \textbf{no constituyen Big Data para Spark}: a esa escala, el costo de serializar, planificar, distribuir y recolectar tareas supera ampliamente el costo del cómputo en sí. El punto de inflexión a partir del cual un motor distribuido compensa su sobrecarga se sitúa por encima del volumen aquí evaluado.

\paragraph{3. Duplicar los nodos \textit{worker} no redujo los tiempos en ninguna de las cuatro herramientas.} Las variaciones entre la Arquitectura A y la B se mantuvieron entre $-2.16$\,\% y $+6.75$\,\%, dentro del ruido de medición, y dos herramientas fueron incluso más lentas en el clúster ampliado. Las causas son distintas según la herramienta y fueron verificadas empíricamente, no supuestas:

\begin{itemize}[leftmargin=*]
    \item \textbf{Polars, Dask y Modin} nunca vieron los nodos adicionales. Polars no es distribuido; Dask levantó un \texttt{LocalCluster} y Modin una instancia local de Ray, ambos confinados al nodo Master por la ausencia de un parámetro \texttt{address} que apuntara al clúster. Los \textit{workers} adicionales permanecieron inactivos durante toda la ejecución.
    \item \textbf{Apache Spark} sí operó sobre YARN, pero tres mecanismos concurrentes impidieron que aprovechara la capacidad extra: la opción \texttt{multiLine=True} forzó la lectura del archivo en una \textbf{única partición}; \textit{Adaptive Query Execution} colapsó las 1,000 particiones de \textit{shuffle} configuradas a solo \textbf{2} particiones reales; y \textit{Dynamic Allocation}, al no detectar tareas encoladas, nunca solicitó más de \textbf{3 executors} sobre 2 de los 4 nodos disponibles. El planificador actuó correctamente: no había trabajo que repartir.
\end{itemize}

\paragraph{4. La integridad del dato se priorizó por encima del rendimiento.} Desactivar \texttt{multiLine=True} habría permitido a Spark particionar la lectura y paralelizarla, pero generaba 500,027 registros en lugar de 500,000, corrompiendo reseñas con saltos de línea internos. Se conservó la opción aun sabiendo que introduce el principal cuello de botella del motor. Un resultado más rápido sobre datos corruptos no es un resultado.

\paragraph{5. La fase de limpieza detectó un defecto del propio conjunto de datos.} El 36.8\,\% de las filas del archivo de trabajo (184,191 de 500,000) resultaron ser copias exactas de otras filas. El origen quedó identificado y verificado: el script de muestreo empleó \texttt{replace=True} sobre un archivo de aproximadamente 496,000 filas, no sobre los 113 millones del dataset completo. El modelo probabilístico del muestreo con reposición predice 316,060 filas distintas para esas condiciones, frente a las 315,809 medidas: una desviación del 0.08\,\%. El hallazgo no invalida los indicadores, porque todos se calculan sobre el conjunto ya deduplicado, pero acota su alcance a la muestra analizada y no a la plataforma completa. Su valor metodológico es que la etapa de validación de calidad cumplió exactamente su propósito: detectar un defecto que el análisis posterior habría arrastrado en silencio.

\paragraph{6. El análisis produjo indicadores de negocio accionables.} La tasa global de recomendación de la muestra es del \textbf{84.43\,\%} (266,640 reseñas positivas sobre 315,809 registros limpios). El \textbf{36.8\,\%} de los registros entregados por la API son duplicados lógicos, por lo que cualquier tablero construido sobre el dato crudo sobreestimaría el volumen de opiniones en más de un tercio. Solo \textbf{609 de 22,992} videojuegos (2.6\,\%) alcanzan las 100 reseñas, lo que obliga a incorporar umbrales de volumen en toda métrica porcentual: sin ese umbral, el ranking por porcentaje de recomendación queda dominado por títulos con una sola reseña. Con el umbral aplicado, \textit{Left 4 Dead 2} (821 reseñas, 100\,\% positivas) encabeza el ranking de fidelización.

\paragraph{7. Dos motores sintácticamente equivalentes no siempre devuelven lo mismo.} La verificación cruzada de las cuatro implementaciones fue exacta en seis de las diez consultas, y reveló una divergencia sistemática en las cuatro que agrupan por videojuego: Polars y Spark reportan 22,992 grupos y Dask y Modin reportan 22,991. La causa es el valor por defecto de un solo parámetro, ya que \texttt{groupby()} de pandas descarta las claves nulas mientras que \texttt{group\_by()} de Polars y \texttt{groupBy()} de Spark las conservan como un grupo. Los 22 registros con \texttt{game} nulo explican la unidad de diferencia. El impacto práctico es nulo (0.007\,\% de los datos, ninguna posición del ranking alterada), pero la lección de ingeniería no lo es: migrar una carga entre motores puede cambiar los resultados sin emitir una sola advertencia, y la defensa es explicitar el criterio en lugar de confiar en el valor por defecto.

\paragraph{8. Hadoop MapReduce confirmó, por una vía independiente, la conclusión sobre el paralelismo.} Los tres programas ejecutados (\texttt{wordmean}, \texttt{secondarysort} y \texttt{terasort}) reprodujeron el mismo patrón observado en Spark: \texttt{secondarysort} (779 KB) y \texttt{wordmean} (2.1 MB) generaron \textbf{un único \textit{input split} y un solo mapper}, mientras que \texttt{terasort} (100 MB) generó \textbf{21 mappers y 7 \textit{reducers}}. El tamaño de bloque de HDFS (128 MB) determina el paralelismo alcanzable, y un clúster de tres nodos no acelera un archivo de 2 MB. Se obtuvo además un resultado transversal sobre depuración distribuida: dos ejecuciones terminaron con estado \texttt{completed successfully} y resultado inválido (\texttt{NaN} por archivo de entrada vacío, \texttt{Infinity} porque \texttt{wordmean} solo lee \texttt{part-r-00000} y sus dos claves se repartieron entre siete \textit{reducers}). En ambos casos fueron los \textbf{contadores del job}, y no los mensajes de estado, los que permitieron el diagnóstico, el mismo principio que resolvió el caso de Spark.

\paragraph{9. Recomendación técnica.} Para el volumen y la topología evaluados, la \textbf{Arquitectura A (1 Master + 2 Workers) con Polars} es la combinación óptima en relación costo-rendimiento, y la Arquitectura B no se justifica económicamente. Escalar el clúster solo aportaría valor bajo alguna de estas condiciones: (a) un volumen de datos un orden de magnitud mayor; (b) una reingeniería del formato de entrada, con conversión previa a Parquet particionado, que eliminaría simultáneamente el cuello de botella de \texttt{multiLine} y habilitaría lectura paralela; o (c) una configuración explícita de Spark con asignación estática de \textit{executors} y AQE desactivado, forzando el paralelismo que el motor descarta por sí mismo. La lección de ingeniería es que \textbf{añadir hardware no acelera un trabajo que el planificador no puede o no necesita dividir}.
